\documentclass[11pt]{article}
\usepackage{mathblatt}
\begin{document}
% Kopfzeile Seite 1 ohne Kennung; die Rückseite (Lösungen, für den Lehrer) bekommt sie nach \begleitteil
\blattkopf{Mathematik}{Basisaufgaben}
\einheitenkopf[][Zettel 15]{Basisaufgaben · Zettel 15}
\zweigzeile{je Aufgabe 1 Punkt · Lösungen auf der Rückseite}
\par\vspace{2pt}\noindent Name \underline{\hspace{7cm}}\hfill Datum \underline{\hspace{3cm}}\hfill Punkte \underline{\hspace{1.2cm}} / 18\par\vspace{6pt}
% \small: volle Seite, kurze Aufgaben zweispaltig (v0.6)
\small

\noindent\begin{minipage}[t]{0.485\linewidth}\raggedright
% 1: Zehnerpotenzschreibweise umwandeln – potenzen-wurzeln-basis-k3-v7 (Original 2025-OS-B1d, ausgedacht: keine Marke)
\begin{aufgabe}{$940\,000 = 9{,}4 \cdot 10^{\square}$ – welche Hochzahl? \leerfeld}
\end{aufgabe}
\end{minipage}\hfill\begin{minipage}[t]{0.485\linewidth}\raggedright
% 2: Uhrzeit aus Startzeit und Dauer berechnen – einheiten-basis-k1-v2 (Original 2021-OS-B1a, ausgedacht: keine Marke)
\begin{aufgabe}{Ein Film beginnt um 13:50 Uhr und dauert $1$ h $25$ min. Wann ist er zu Ende? \leerfeld Uhr}
\end{aufgabe}
\end{minipage}\par

\noindent\begin{minipage}[t]{0.485\linewidth}\raggedright
% 3: Größen vergleichen – einheiten-basis-k2-v3 (Original 2026-FOR-B1d, ausgedacht: keine Marke)
\begin{aufgabe}{Vergleiche $0{,}3$ l und $30$ ml. Setze das richtige Zeichen ein: $<$, $=$ oder $>$. $0{,}3$ l \leerfeld $30$ ml}
\end{aufgabe}
\end{minipage}\hfill\begin{minipage}[t]{0.485\linewidth}\raggedright
% 4: Strecke aus Teilstrecken berechnen – einheiten-basis-k4-v1 (Original 2018-OS-B1b, ausgedacht: keine Marke)
\begin{aufgabe}{Ein Brett ist $3$ m lang. Davon werden $4$ Stücke von je $45$ cm abgesägt. Wie lang ist der Rest? \leerfeld[cm]}
\end{aufgabe}
\end{minipage}\par

\noindent\begin{minipage}[t]{0.485\linewidth}\raggedright
% 5: Proportionale Zuordnung Dreisatz – zuordnungen-basis-k2-v3 (Original 2023-OS-B1a, ausgedacht: keine Marke)
\begin{aufgabe}{Ein Drucker druckt $12$ Seiten in $2$ min. Wie viele Seiten druckt er in $5$ min? \leerfeld Seiten}
\end{aufgabe}
\end{minipage}\hfill\begin{minipage}[t]{0.485\linewidth}\raggedright
% 6: Term durch Zusammenfassen gleichartiger Glieder vereinfachen – terme-basis-k1-v2 (Original 2025-GYM-B2a, ausgedacht: keine Marke)
\begin{aufgabe}{Vereinfache den Term $9y - y^2 - 4y$ und berechne seinen Wert für $y = 3$.}
\end{aufgabe}
\end{minipage}\par

\noindent\begin{minipage}[t]{0.485\linewidth}\raggedright
% 7: Lineare Gleichung lösen – lineare-gleichungen-basis-k1-v9 (Original 2024-OS-B1d, ausgedacht: keine Marke)
\begin{aufgabe}{Löse die Gleichung $5 \cdot (x + 3{,}5) = 0$. x = \leerfeld}
\end{aufgabe}
\end{minipage}\hfill\begin{minipage}[t]{0.485\linewidth}\raggedright
% 8: Lösung durch Einsetzen prüfen – quadratische-gleichungen-basis-k1-v7 (Original 2025-OS-B1h, ausgedacht: keine Marke)
\begin{aufgabe}{Welcher Wert erfüllt $x \cdot (x - 5) = 6$? \\ \kreuz{$x = 1$} \kreuz{$x = -6$} \kreuz{$x = 6$} \kreuz{$x = 5$}}
\end{aufgabe}
\end{minipage}\par

\noindent\begin{minipage}[t]{0.485\linewidth}\raggedright
% 9: Quadratseite aus Fläche berechnen – flaechen-basis-k2-v2 (Original 2020-OS-B1d, ausgedacht: keine Marke)
\begin{aufgabe}{Ein quadratisches Beet hat den Flächeninhalt $81$ m². Wie lang ist eine Seite? \leerfeld[m]}
\end{aufgabe}
\end{minipage}\hfill\begin{minipage}[t]{0.485\linewidth}\raggedright
% 10: Umfang Rechteck berechnen – flaechen-basis-k5-v1 (Original 2026-FOR-B1h, ausgedacht: keine Marke)
\begin{aufgabe}{Ein Beet ist rechteckig, $4{,}5$ m lang und $2$ m breit. Berechne den Umfang. u = \leerfeld[m]}
\end{aufgabe}
\end{minipage}\par

\noindent\begin{minipage}[t]{0.485\linewidth}\raggedright
% 11: Volumen Würfel berechnen – koerper-basis-k3-v2 (Original 2026-FOR-B1f, ausgedacht: keine Marke)
\begin{aufgabe}{Ein Würfel hat die Kantenlänge $5$ cm. Gib sein Volumen an. \leerfeld[cm³]}
\end{aufgabe}
\end{minipage}\hfill\begin{minipage}[t]{0.485\linewidth}\raggedright
% 12: Trigonometrische Gleichung nach Seite umstellen – trigonometrie-basis-k1-v2 (Original 2020-OS-B1j, ausgedacht: keine Marke)
\begin{aufgabe}{$\mathrm{cos}\,60^\circ = \frac{9}{x}$ – berechne $x$. x = \leerfeld}
\end{aufgabe}
\end{minipage}\par

% 13: Vorzeichenregel anwenden – rationale-zahlen-basis-k2-v2 (Original 2015-OS-B1g, ausgedacht: keine Marke)
\begin{aufgabe}{Eine positive Zahl wird durch eine negative Zahl geteilt. Das Ergebnis ist … \\ \kreuz{immer positiv} \kreuz{immer negativ} \kreuz{Das kann man nicht entscheiden.}}
\end{aufgabe}

% 14: Prozent und Anteil umwandeln – prozentrechnung-basis-k2-v4 (Original 2022-OS-B1f, ausgedacht: keine Marke)
\begin{aufgabe}{Jede fünfzigste Schraube ist fehlerhaft – wie viel Prozent? Kreuze an. \\ \kreuz{$2\,\%$} \kreuz{$5\,\%$} \kreuz{$20\,\%$} \kreuz{$50\,\%$}}
\end{aufgabe}

% 15: Eigenschaft einer Figur zuordnen – winkel-dreiecke-basis-k5-v3 (Original 2026-FOR-B1c, ausgedacht: keine Marke)
\begin{aufgabe}{Kreuze die richtige Ergänzung an: „In jedem Quadrat …“\\ \kreuz{… gibt es genau eine Symmetrieachse.}\\ \kreuz{… sind zwei Winkel stumpf.}\\ \kreuz{… sind die Diagonalen gleich lang und stehen senkrecht aufeinander.}}
\end{aufgabe}

% 16: Wahrscheinlichkeit einstufig – bruchrechnung-basis-k1-v6 (Original 2016-OS-B1f, ausgedacht: keine Marke)
\begin{aufgabe}{Ein Glücksrad hat 10 gleich große Felder mit den Zahlen 1 bis 10. Wie groß ist die Wahrscheinlichkeit für eine durch 3 teilbare Zahl? \leerfeld}
\end{aufgabe}

% 17: Bruchteil einer Fläche bestimmen – brueche-dezimalzahlen-basis-k1-v12 (Original 2026-FOR-B1b, ausgedacht: keine Marke)
\begin{aufgabe}{\begin{minipage}[t]{0.6\linewidth}Eine Pizza ist in gleich große Stücke geschnitten; die grauen Stücke sind schon gegessen. Welcher Anteil der Fläche ist grau? Gib ihn als gekürzten Bruch an. \leerfeld\end{minipage}\hfill\begin{minipage}[t]{0.37\linewidth}\vspace{0pt}
\bruchkreis[0.9]{6}{8}
\end{minipage}}
\end{aufgabe}

% 18: Graph nach Eigenschaft auswählen – lineare-funktionen-basis-k2-v2 (Original 2016-OS-B1c, ausgedacht: keine Marke)
\begin{aufgabe}{\begin{minipage}[t]{0.6\linewidth}Welche Gerade fällt? Kreuze an. \\ \kreuz{f} \kreuz{g}\end{minipage}\hfill\begin{minipage}[t]{0.37\linewidth}\vspace{0pt}
\begin{ksys}[xmin=-5,xmax=5,ymin=-5,ymax=5,klein] \gerade{1}{2}{f} \gerade{-1.5}{-1}{g} \end{ksys}
\end{minipage}}
\end{aufgabe}

\begleitteil
\blattkopf{Basisaufgaben · Zettel 15}{BAS-Z15 · Lösungen}
\erg{1}{$5$ – das Komma wandert fünf Stellen nach links \hfill {\footnotesize\textit{Zehnerpotenzschreibweise umwandeln · 2025}}}
\erg{2}{13:50 Uhr $+ 1$ h $=$ 14:50 Uhr, $+ 25$ min $=$ 15:15 Uhr \hfill {\footnotesize\textit{Uhrzeit aus Startzeit und Dauer berechnen · 2021}}}
\erg{3}{$0{,}3$ l $= 300$ ml, also $0{,}3$ l $>$ $30$ ml \hfill {\footnotesize\textit{Größen vergleichen · 2026}}}
\erg{4}{$3$ m $= 300$ cm; $300 - 4 \cdot 45 = 300 - 180 = 120$ cm \hfill {\footnotesize\textit{Strecke aus Teilstrecken berechnen · 2018}}}
\erg{5}{Für $1$: $12 : 2 = 6$; $6 \cdot 5 = 30$ Seiten \hfill {\footnotesize\textit{Proportionale Zuordnung Dreisatz · 2023}}}
\erg{6}{$9y - 4y = 5y$, Term $5y - y^2$; für $y = 3$: $5 \cdot 3 - 3^2 = 15 - 9 = 6$ \hfill {\footnotesize\textit{Term durch Zusammenfassen gleichartiger Glieder vereinfachen · 2025}}}
\erg{7}{$x + 3{,}5 = 0$, also $x = -3{,}5$ \hfill {\footnotesize\textit{Lineare Gleichung lösen · 2024}}}
\erg{8}{$x = 6$, denn $6 \cdot (6 - 5) = 6 \cdot 1 = 6$ \hfill {\footnotesize\textit{Lösung durch Einsetzen prüfen · 2025}}}
\erg{9}{$9 \cdot 9 = 81$, also $9$ m \hfill {\footnotesize\textit{Quadratseite aus Fläche berechnen · 2020}}}
\erg{10}{$u = 2 \cdot (4{,}5 + 2) = 13$ m \hfill {\footnotesize\textit{Umfang Rechteck berechnen · 2026}}}
\erg{11}{$5 \cdot 5 \cdot 5 = 125$ cm³ (nicht $5 \cdot 5 = 25$) \hfill {\footnotesize\textit{Volumen Würfel berechnen · 2026}}}
\erg{12}{$x = 9 : \mathrm{cos}\,60^\circ = 9 : 0{,}5 = 18$ \hfill {\footnotesize\textit{Trigonometrische Gleichung nach Seite umstellen · 2020}}}
\erg{13}{immer negativ \hfill {\footnotesize\textit{Vorzeichenregel anwenden · 2015}}}
\erg{14}{$2\,\%$ ($\frac{1}{50} = \frac{2}{100}$) \hfill {\footnotesize\textit{Prozent und Anteil umwandeln · 2022}}}
\erg{15}{… sind die Diagonalen gleich lang und stehen senkrecht aufeinander. \hfill {\footnotesize\textit{Eigenschaft einer Figur zuordnen · 2026}}}
\erg{16}{$\frac{3}{10}$ (3, 6, 9) \hfill {\footnotesize\textit{Wahrscheinlichkeit einstufig · 2016}}}
\erg{17}{$\frac{6}{8} = \frac{3}{4}$ ($6$ von $8$ gleichen Teilen) \hfill {\footnotesize\textit{Bruchteil einer Fläche bestimmen · 2026}}}
\erg{18}{g (von links nach rechts gelesen geht sie nach unten) \hfill {\footnotesize\textit{Graph nach Eigenschaft auswählen · 2016}}}
\end{document}
